\documentclass[12pt]{article}
\usepackage[a4paper,margin=2cm]{geometry}
\usepackage{amsmath,amssymb}
\usepackage{xepersian}
\settextfont[
  Path=../1401-2/HW2 - official solution/,
  Extension=.ttf,
  BoldFont=*Bd,
  ItalicFont=*It,
  BoldItalicFont=*BdIt
]{XB Niloofar}
\setdigitfont[
  Path=../1401-2/HW2 - official solution/,
  Extension=.ttf,
  BoldFont=*Bd,
  ItalicFont=*It,
  BoldItalicFont=*BdIt
]{XB Niloofar}
\setlength{\parindent}{0pt}
\setlength{\parskip}{0.55em}

\begin{document}
\begin{center}
  {\Large\textbf{اصلاحیهٔ پاسخ رسمی تمرین سوم ـ نیم‌سال ۱۴۰۲--۱}}
\end{center}

این برگه فقط موارد نادرست پاسخ منتشرشده را اصلاح می‌کند و باید همراه فایل اصلی خوانده شود.

\section*{سؤال ۱}
در قاعدهٔ مرکب سیمپسون، تعداد زیربازه‌ها باید زوج باشد. با اعمال کران خطای سیمپسون و گردکردن رو به بالای تعداد زیربازه‌ها به نزدیک‌ترین عدد زوج، مقادیر درست عبارت‌اند از
\[
  n_{\text{الف}}=52,\qquad
  n_{\text{ب}}=40,\qquad
  n_{\text{ج}}=42.
\]
در بخش الف، بیشینه‌ی دقیق
\(\lvert4\sin x+x\cos x\rvert\)
روی
\([0,\pi]\)
برابر تقریباً
\(4.20256854\)
است و کمترین تعداد زوج را
\(52\)
می‌دهد؛ مقدار ۶۰ فقط از کران محافظه‌کارانه‌ی
\(4+\pi\)
به‌دست می‌آید. در بخش‌های ب و ج، کران‌های محاسبه‌شده‌ی پاسخ اصلی
\(38.2177\)
و
\(40.5360\)
به ۳۹ و ۴۱ گرد شده‌اند؛ چون سیمپسون تعداد زیربازه‌ی زوج می‌خواهد، مقادیر درست ۴۰ و ۴۲ هستند.

\section*{سؤال ۴}
ضرایب درست فرمول پنج‌نقطه‌ای خواسته‌شده برابرند با
\[
 f'(x_0)\simeq
 \frac{
 -\frac14 f(x_0-h)
 -\frac56 f(x_0)
 +\frac32 f(x_0+h)
 -\frac12 f(x_0+2h)
 +\frac1{12}f(x_0+3h)}
 {h}.
\]
برای \(f(x)=\sin x\cos x\)، \(x_0=1\) و \(h=0.1\) نتیجهٔ عددی
\[
  f'(1)\simeq-0.4161915015
\]
است؛ مقدار دقیق \(\cos(2)\simeq-0.4161468365\) است.

\section*{سؤال ۵}
برای تفاضل مرکزی
\[
  D_hf(x)=\frac{\widetilde f(x+h)-\widetilde f(x-h)}{2h},
\qquad
 |\widetilde f-f|\leq\alpha,
\]
کران خطای ناشی از گردکردن برابر \(\alpha/h\) است، نه \(\alpha/(2h)\). در نتیجه
\[
  |f'(x)-D_hf(x)|
  \leq \frac{\alpha}{h}+\frac{M h^2}{6},
  \qquad
  M\geq |f^{(3)}(\xi)|
  \quad\text{روی یک بازه‌ی ثابت شامل همه‌ی نقاط مجاز }x\pm h.
\]
کمینه‌کردن این کران می‌دهد
\[
  h_{\mathrm{opt}}=\left(\frac{3\alpha}{M}\right)^{1/3}.
\]
گام بسیار کوچک خطای گردکردن را غالب می‌کند و گام بسیار بزرگ خطای برشی را افزایش می‌دهد.
\end{document}
